\pagebreak

\chapter{Specific Limitations}

    \section{multiply, divide and remainder operators}
    
        \subsection{Virtex-E}

            In the Virtex-E there are no hardware multiply or
            divide functions. The multiply, divide and remainder operators
            carry out these operations using cascaded adders or subtractors.
            Since these are expression operators they cannot be pipelined. The
            result is that these operations have a large data delay, typically
            30ns to 40ns.
            
            \index{constraint!timing}
            \index{timing constraint}
            If the operation can be allowed to take place over several
            clock cycles between static variables then the timing constraints
            can be changed to suit, e.g.
            \begin{lstlisting}[gobble=12, language=threepl]
               static(a, "uint:8"); attribute(a, {tnm="SRC"});
               static(b, "uint:8"); attribute(b, {tnm="SRC"});
               static(c, "uint:16"); attribute(c, {tnm="DEST"});
               maxdelay("SRC", "DEST", 40.0e-9); // 40ns path constraint

               c <- a * b;     // slow operation!
            \end{lstlisting}
            
            If the operation must be at clock rate then pipelined 3PL library
            modules should be used for multiply, divide and remainder.
    
        \subsection{Virtex-II or Virtex-II Pro}

            In the Virtex-II or Virtex-II Pro there are hardware multipliers
            but no hardware dividers.

            The multiply operator uses one hardware multiplier for operands up
            to 18 bits if signed, 17 bits if unsigned. If one operand is between
            19 and 35 bits signed, or 18 to 34 bits unsigned, then two hardware
            multipliers plus an adder are used. If both operands are between 19
            and 35 bits signed, or 18 to 34 bits unsigned, then four hardware
            multipliers plus two adders are used.
            
            No attempt has been made to share hardware multipliers between
            products of small operands. Hardware multipliers are a limited
            resource and care must be taken not to exceed the number available
            in the particular FPGA being used.
                       
            There are no hardware dividers and divide and remainder operators
            carry out these operations using cascaded adders and subtractors.
            Since these are expression operators they cannot be pipelined. The
            result is that these operations have a large data delay, typically
            30ns to 40ns. For approaches to this see the previous subsection
            above (Virtex-E).
    
        \subsection{Virtex-5}

            In the Virtex-5 there are hardware multipliers but no hardware
            dividers. The multipliers are contained in the DSP48E blocks.

            The multiply operator uses one hardware multiplier for operands up
            to 25 bits by 18 bits signed. Unsigned operands have a 0 sign bit
            appended at the most significant end and thus become one bit
            larger.
            
            Selectvalue, static and queue mode variables which are assigned
            arithmetic or logical expressions can employ DSP blocks implicitly
            by setting the DSP attribute. In the case of static variables this
            will also use the DSP block output register rather than using
            flip-flops (but not if the variable is initialised to a value other
            than zero since the DSP block output register cannot be
            initialised). Selectvalue, static and queue mode variables can
            have the ALU attribute (and not the DSP attribute) in which case
            additions and subtractions in assigned expressions will be
            combined in a common ALU implemented in CLBs.
            
            To take advantage of the detailed architecture of the DSP blocks
            they can be instantiated directly by calling
            {\bf DSP48E()} \autorefph{sec:lpdsp48e} in library xc5v.3pl.
            Alternatively there are some useful modules providing a simpler
            interface in that library - \\
            adder/subtracter - {\bf dspaddsub} \autorefph{sec:lmdspaddsub} \\
            counter - {\bf dspcounter} \autorefph{sec:lmdspcounter} \\
            logical functions - {\bf dsplog} \autorefph{sec:lmdsplog} \\
            multiplier - {\bf dspmuladd} \autorefph{sec:lmdspmul} \\
            multiplier adder - {\bf dspmuladd} \autorefph{sec:lmdspmuladd} \\
            multiplier subtracter - {\bf dspmulsub} \autorefph{sec:lmdspmulsub} \\
            subtracter multiplier - {\bf dspsubmul} \autorefph{sec:lmdspsubmul} \\
            
            If these modules are used with no DSP48E registers then the
            blocks can be used as a combinatorial elements in simple
            static and queue variable assignments. If DSP48E internal
            pipeline registers are used then the pipeline delays must
            be taken into account. This complicates the programming style
            considerably and is a tradeoff for the performance and resource
            usage advantages gained.

    
        \subsection{Virtex-6}
        
            At the moment the support is identical to that for virtex-5 -
            see above.




    \section{registered output memory}
    
        Registered output memory in Xilinx FPGAs is block RAM. These are a
        limited resource.
    
        \subsection{XCV300E}
        
            Block RAMs in the XCV3 have 4,096 bits which
            can be configured as -
            
            \begin{tabular}{l l l}
            \hline
            data width & address width & depth \\
            \hline
            1  & 12 & 4096 \\
            2  & 11 & 2048 \\
            4  & 10 & 1024 \\
            8  & 9  & 512  \\
            16 & 8  & 256  \\
            \hline
            \end{tabular}

            i.e. anywhere from 8 bit address to 12 bit address.
            
            The 3PL compiler will select blocks of appropriate size
            and configuration to provide depths of 256 to 4096, using as
            many blocks as required to match the data width.
    
        \subsection{XCV3, XC2VP}
        
            Block RAMs in the XCV3 have 8,432 bits which
            can be configured as -
            
            \begin{tabular}{l l l}
            \hline
            data width & address width & depth \\
            \hline
            1  & 14 & 16384 \\
            2  & 13 & 8192  \\
            4  & 12 & 4096  \\
            9  & 11 & 2048  \\
            18 & 10 & 1024  \\
            36 & 9  & 512   \\
            \hline
            \end{tabular}

            i.e. anywhere from 9 bit address to 14 bit address
            (the unusual data widths result from the use of extra 'parity'
            bits at the wider configurations).
            
            The 3PL compiler will select blocks of appropriate size
            and configuration to provide depths of 512 to 16384, using as
            many blocks as required to match the data width.
    
        \subsection{Virtex-5, Virtex-6}
        
            Block RAM elements are grouped in pairs, each block containing
            18Kbits. The pair can be configured as a single 36Kbit block.
            18K blocks can be configured as - 
            
            \begin{tabular}{l l l}
            data width & address width & depth \\
            \hline
            1  & 14 & 16384 \\
            2  & 13 & 8192 \\
            4  & 12 & 4096 \\
            9  & 11 & 2048 \\
            18 & 10 & 1024 \\
            \end{tabular}

            36K blocks can be configured as - 
            
            \begin{tabular}{l l l}
            data width & address width & depth \\
            \hline
            1  & 15 & 32768 \\
            2  & 14 & 16384 \\
            4  & 13 & 8192  \\
            9  & 12 & 4096  \\
            18 & 11 & 2048  \\
            36 & 10 & 1024  \\
            \end{tabular}
            
            The 3PL compiler will select blocks of appropriate size and
            configuration to provide depths of 1024 to 32768, using as many
            blocks as required to match the data width.

            Block RAM is also used to implement deep queues. For queues whose
            depth is 512, 1024, 2048, 4096 or 8192 a FIFO primitive element is
            used. This uses a block RAM with inbuilt FIFO control logic, saving
            considerable external CLB logic and probably giving improved
            performance (but note that the clock to output time of nearly 2nS
            still persists). The hardware FIFO is not used in the case of
            asynchronous queues on which the queuereadstatus() and
            queuewritestatus() have been used.

    \section{selectvalue variables}

        Selectvalue variables implement buses using CLB logic unless explicitly
        directed to use 3-state buffers driving partitionable 'long lines'. This
        is done by setting attribute "threestate" on individual selectvalue
        variables. These are only available on older FPGAs such as
        Virtex-E, Virtex-II and Virtex-II Pro.
        

    \section{high speed serial}
    
    
        \subsection{Virtex-II Pro}
        
            MGT high speed serial ports are interfaced using library {\bf
            gt\_custom\_16()} \autorefph{sec:lmgtcustom16} in library {\bf
            xilinx/mgt.3pl}.
    
        \subsection{Virtex-5, Virtex-6}
        
            GTP high speed serial ports are interfaced using library
            module {\bf gtp\_dual\_16()} \autorefph{sec:lmgtpdual16} in
            library {\bf xilinx/gtp\_dual.3pl}.
            
            GTX high speed serial ports are not yet supported.
